\chapter{Caracteres y transformada de Gelfand}\label{cap:v3-12}

Un carácter convierte cada elemento de un álgebra conmutativa en un número, respetando sumas, productos e identidad. Reunir todos esos modos de evaluación produce un espacio compacto; cada elemento del álgebra se convierte entonces en una función continua sobre ese espacio. La construcción no siempre conserva toda la información: su núcleo medirá precisamente lo que los caracteres no pueden detectar.

\section{Caracteres e ideales maximales}\label{lec:v3-24}
\fuentelectura{24}{La transformada de Gelfand}

En este capítulo $A$ es un álgebra de Banach compleja, conmutativa, unitaria y no nula, con $\norm1=1$.

\begin{definicion}
Un \emph{carácter} de $A$ es un homomorfismo complejo lineal y multiplicativo $\varphi:A\to\C$ que no es idénticamente cero. El conjunto de caracteres se denota $\Delta(A)$.
\end{definicion}
\index{carácter}\index{espacio de ideales maximales}

La exclusión del homomorfismo cero es indispensable. Si $\varphi(a)\ne0$, la identidad $\varphi(1)\varphi(a)=\varphi(a)$ da $\varphi(1)=1$. No necesitamos suponer continuidad: será consecuencia de la estructura algebraica.

\begin{proposicion}\label{v3:caracteres-maximales}
Cada carácter satisface $\varphi(a)\in\sigma(a)$ y $\norm\varphi=1$. Su núcleo es un ideal maximal cerrado. Recíprocamente, cada ideal maximal $M$ es el núcleo de un único carácter, determinado por
\[
 \varphi_M(a)=\lambda\quad\Longleftrightarrow\quad a-\lambda1\in M.
\]
Además,
\[
 \sigma(a)=\{\varphi(a):\varphi\in\Delta(A)\}.
\]
\end{proposicion}
\begin{proof}
Si $a-\varphi(a)1$ fuera invertible, aplicar $\varphi$ a su producto con la inversa daría $0=1$. Así $|\varphi(a)|\le r(a)\le\norm a$, y $\norm\varphi\le1$; evaluar en $1$ da la igualdad. El núcleo es cerrado por continuidad y es ideal por multiplicatividad. Como $\varphi$ es sobreyectiva ($\varphi(\lambda1)=\lambda$), el cociente por su núcleo es $\C$ y el núcleo es maximal.

Para la recíproca, un ideal maximal de un álgebra de Banach unitaria es cerrado, y $A/M$ es un álgebra de Banach de división. Gelfand--Mazur identifica cada clase con un único $\lambda[1]$. La unicidad se sigue también de que $\lambda1\in M$ con $\lambda\ne0$ implicaría $1\in M$. Sumar y multiplicar clases prueba que $\varphi_M$ es lineal, multiplicativa y unital; su núcleo es exactamente $M$. Cualquier carácter con ese núcleo debe asignar a $a$ el mismo escalar, lo cual prueba la unicidad.

Ya se probó que todos los valores de caracteres pertenecen al espectro. Si $\lambda\in\sigma(a)$, el ideal $(a-\lambda1)A$ es propio, pues de otro modo $a-\lambda1$ tendría inversa en el álgebra conmutativa. Está contenido en un ideal maximal $M$ y entonces $\varphi_M(a)=\lambda$. Esto prueba la inclusión restante y, en particular, la existencia de caracteres.
\end{proof}

\begin{ejemplo}[Series absolutamente convergentes]
En $\ell^1(\N_0)$ con el producto de Cauchy, para $|z|\le1$ defina
\[
 \varphi_z(a)=\sum_{n=0}^{\infty}a_nz^n.
\]
La serie converge absolutamente y $|\varphi_z(a)|\le\norm a_1$, por lo que la aplicación es lineal continua. Además $\varphi_z(e_0)=1$. La doble serie $\sum_{j,k}|a_jb_kz^{j+k}|$ es finita; reagruparla por $n=j+k$ demuestra $\varphi_z(a*b)=\varphi_z(a)\varphi_z(b)$. Más adelante verificaremos que éstos son todos los caracteres.
\end{ejemplo}

\section{La topología del espacio de caracteres}

Dotamos a $\Delta(A)\subset A^*$ de la topología débil estrella: una red $\varphi_\alpha$ converge a $\varphi$ si y sólo si $\varphi_\alpha(a)\to\varphi(a)$ para cada $a\in A$. Esta topología registra la convergencia de todas las evaluaciones, sin exigir convergencia en norma de los funcionales.

\begin{teorema}\label{v3:compacto-caracteres}
El espacio $\Delta(A)$ es compacto Hausdorff. Si $A$ es separable, es metrizable.
\end{teorema}
\begin{proof}
La topología débil estrella es Hausdorff porque dos funcionales distintos difieren en algún elemento de $A$. La bola dual cerrada unitaria es compacta por Banach--Alaoglu. Dentro de ella,
\[
 \Delta(A)=\{\varphi:\varphi(1)=1,\ \varphi(ab)=\varphi(a)\varphi(b)\text{ para todo }a,b\in A\}
\]
es cerrado: cada igualdad se conserva bajo convergencia puntual de funcionales. En particular, la ecuación $\varphi(1)=1$ impide que una red de caracteres tenga límite cero; verificar únicamente la multiplicatividad no bastaría.

Si $\{a_n\}$ es denso en la bola unitaria de $A$, la fórmula
\[
 d(\varphi,\psi)=\sum_{n=1}^{\infty}2^{-n}|\varphi(a_n)-\psi(a_n)|
\]
define una métrica en la bola dual: la serie converge pues cada diferencia es a lo sumo $2$, y la densidad garantiza que distancia cero implica igualdad. La convergencia débil estrella implica convergencia para esta métrica al controlar primero una suma finita y luego la cola geométrica. Recíprocamente, la convergencia de las evaluaciones en $a_n$, junto con la cota uniforme de las normas, permite aproximar cualquier $a$ por un múltiplo de algún $a_n$ y obtener convergencia en $a$. Así la métrica induce exactamente la topología deseada en $\Delta(A)$.
\end{proof}

\section{La representación por funciones continuas}

\begin{definicion}
Para $a\in A$, su \emph{transformada de Gelfand} es la función
\[
 \widehat a:\Delta(A)\to\C,\qquad \widehat a(\varphi)=\varphi(a).
\]
La aplicación de Gelfand es $\mathcal G:A\to C(\Delta(A))$, $a\mapsto\widehat a$.
\end{definicion}
\index{transformada de Gelfand}

\begin{teorema}\label{v3:transformada-gelfand}
La aplicación $\mathcal G$ es un homomorfismo unital continuo, y
\[
 \widehat a(\Delta(A))=\sigma(a),\qquad
 \norm{\widehat a}_\infty=r(a)\le\norm a.
\]
Si $f$ es holomorfa cerca de $\sigma(a)$, entonces
\[
 \widehat{f(a)}=f\circ\widehat a.
\]
\end{teorema}
\begin{proof}
La función $\widehat a$ es continua por la definición de la topología débil estrella: es precisamente una de sus evaluaciones generadoras. La desigualdad $|\varphi(a)-\psi(a)|\le\norm{\varphi-\psi}\norm a$ sólo probaría continuidad para la topología de la norma, que es más fuerte y no resuelve esta cuestión. La linealidad y la multiplicatividad de cada carácter dan, punto por punto,
$\widehat{\alpha a+\beta b}=\alpha\widehat a+\beta\widehat b$, $\widehat{ab}=\widehat a\,\widehat b$ y $\widehat1=1$. La descripción del espectro de la proposición~\ref{v3:caracteres-maximales} da las igualdades de imagen y norma, de donde resulta continuidad de $\mathcal G$.

Finalmente, un carácter conmuta con la integral vectorial por ser lineal continuo y envía $(z1-a)^{-1}$ a $(z-\varphi(a))^{-1}$. La fórmula de Cauchy implica
\[
 \varphi(f(a))=\frac1{2\pi i}\int_\Gamma
 \frac{f(z)}{z-\varphi(a)}\,dz=f(\varphi(a)).
\]
Esto vale para cada carácter, que es exactamente la última identidad.
\end{proof}

\section{Ejercicios sobre el radical y la representación}

\begin{ejercicio}\label{ej:v3-l24-1}
Un elemento $a$ es \emph{cuasinilpotente} si $\sigma(a)=\{0\}$. (i) Pruebe que
\[
 J=\{a\in A:ab\text{ es cuasinilpotente para todo }b\in A\}
\]
es un ideal cerrado, llamado radical de Jacobson, y que es la intersección de todos los ideales maximales. (ii) Pruebe que todo elemento cuasinilpotente es divisor topológico de cero. (iii) Pruebe $\ker\mathcal G=J$. Verifique además que cuasinilpotente no implica nilpotente para el operador de desplazamiento ponderado
\[
 T(x_1,x_2,\ldots)=(0,x_1,x_2/2,x_3/3,\ldots)
 \quad\text{en }\ell^\infty.
\]
\end{ejercicio}


\begin{ejercicio}\label{ej:v3-l24-2}
Pruebe que todo carácter de $\ell^1(\N_0)$ con convolución de Cauchy es $\varphi_z$ para un único $z$ con $|z|\le1$.
\end{ejercicio}


\begin{ejercicio}\label{ej:v3-l24-3}
Demuestre que todo carácter es una aplicación abierta.
\end{ejercicio}


\begin{ejercicio}\label{ej:v3-l24-4}
Pruebe $\sigma(a+b)\subset\sigma(a)+\sigma(b)$ y $\sigma(ab)\subset\sigma(a)\sigma(b)$, y dé ejemplos en que ambas inclusiones sean propias.
\end{ejercicio}


\begin{ejercicio}\label{ej:v3-l24-5}
Pruebe que $a$ es invertible si y sólo si $\widehat a$ es invertible en $C(\Delta(A))$.
\end{ejercicio}


\begin{ejercicio}\label{ej:v3-l24-6}
Pruebe $r(a)=\norm{\widehat a}_\infty$. Deduzca que $\mathcal G$ es isométrica si y sólo si $\norm{a^2}=\norm a^2$ para cada $a\in A$.
\end{ejercicio}


El radical explica por qué una representación por funciones escalares puede perder información. La identidad $C^*$ eliminará esa pérdida en el caso conmutativo y permitirá identificar el álgebra completa con todas las funciones continuas de su espacio de caracteres. La exposición se complementa con \cite{maxgelfand,murphy1990,rudin1991}.
